\chapter{Hypothesis tests and confidence sets}
\label{sec:p15-hypothesis-tests-confidence}

A test controls a probability of rejection under a specified sampling
model. A confidence procedure controls the probability that its random
reported set contains a fixed parameter. This chapter supplies the
simple-hypothesis comparison of Neyman and Pearson
\cite[pp.~293--301]{NeymanPearson1933Tests} and Neyman's acceptance-region
construction \cite[pp.~347--355]{Neyman1937Confidence}, with explicit
measurability, randomization and discrete boundary conditions.
The binomial applications are derived from these constructions.

\begin{hbdefinition}{def:p15-statistical-experiment}{Experiments, tests, errors and power}
A statistical experiment consists of a measurable sample space
\((\mathsf X,\mathscr X)\), a parameter set \(\Theta\), and a
probability measure \(P_\theta\) on that space for each \(\theta\).
The observed data \(X\) have the specified law \(P_\theta\).
A simple hypothesis specifies one law; a composite hypothesis
allows several laws. Let \(\Theta_0\ne\varnothing\) be the null
parameter set and \(\Theta_1\) a disjoint alternative set.

A randomized test is a measurable function
\(\varphi:\mathsf X\to[0,1]\). Its value is the probability of
rejecting the null at that observation. A nonrandomized test is
an indicator \(\mathbf1_R\) of a measurable rejection region.
Define
\[
 \beta_\varphi(\theta)=E_\theta\varphi(X),\qquad
 \operatorname{size}(\varphi)=\sup_{\theta\in\Theta_0}\beta_\varphi(\theta).
\]
The test has level \(\alpha\in[0,1]\) if its size is at most
\(\alpha\). Its type-I error probability at a null law is
\(\beta_\varphi(\theta)\); its power at an alternative is that
same rejection probability, and its type-II error probability
is \(1-\beta_\varphi(\theta)\).
A most-powerful level-\(\alpha\) test maximizes power against
the specified simple alternative. A uniformly most-powerful
test does so against every member of the stated alternative family.
Nonrejection is a decision outcome; it does not assert that
the null hypothesis has been proved.
\end{hbdefinition}

\begin{hblemma}{thm:p15-test-randomization}{proved=annex:p15.test.randomization}{The exact product-space implementation}
Let \(U\) be uniform on \((0,1)\), independent of \(X\).
Write \(\lambda_{01}\) for this uniform probability measure.
Rejecting when \(U\leq\varphi(X)\) has probability
\(E_\theta\varphi(X)\) under every \(P_\theta\).
The rejection event is measurable. Both \(\varphi=0\) and
\(\varphi=1\) have their stated deterministic meanings.
\end{hblemma}

\section{The two-density comparison}

\begin{hbdefinition}{def:p15-likelihood-test}{Densities and the likelihood ordering}
Let \(P_0,P_1\) have nonnegative measurable densities \(f_0,f_1\)
relative to a common measure \(\mu\), with integrals 1.
Choose versions finite everywhere by changing any infinite-value
null sets to zero. Define the extended likelihood ratio
\[
 L(x)=
 \begin{cases}
 f_1(x)/f_0(x),&f_0(x)>0,\\
 +\infty,&f_0(x)=0<f_1(x),\\
 0,&f_0(x)=f_1(x)=0.
 \end{cases}
\]
For a finite threshold \(c\geq0\), write
\[
 \varphi_{c,\gamma}
   =\mathbf1_{\{f_1>cf_0\}}
        +\gamma\mathbf1_{\{f_1=cf_0,\ f_0>0\}},
                         \qquad0\leq\gamma\leq1.
\]
The common zero-density set has no effect under either law.
Points possible only under \(P_1\) are rejected without
spending any null rejection probability.
\end{hbdefinition}

\begin{hblemma}{thm:p15-likelihood-threshold}{proved=annex:p15.likelihood.threshold}{Existence of a threshold with exact prescribed size}
For \(0<\alpha<1\), there is a finite \(c\geq0\) such that
\[
                      P_0(L>c)\leq\alpha\leq P_0(L\geq c).
\]
If \(q=P_0(L=c)>0\), choose
\(\gamma=(\alpha-P_0(L>c))/q\); if \(q=0\), choose
\(\gamma=0\). Then \(E_0\varphi_{c,\gamma}=\alpha\).
This includes a threshold atom and \(c=0\).
At the size endpoints use
\[
                  \varphi_0=\mathbf1_{\{f_0=0\}},\qquad
                  \varphi_1=1.
\]
\end{hblemma}

\begin{hbtheorem}{thm:p15-neyman-pearson}{proved=annex:p15.neyman.pearson}{The dominated Neyman--Pearson lemma with the complete converse}
Let \(0<\alpha<1\), and let \(\varphi^*=\varphi_{c,\gamma}\)
be the preceding exact-size test. For every test \(\psi\)
with \(E_0\psi\leq\alpha\),
\[
 E_1\varphi^*-E_1\psi
 =\int(\varphi^*-\psi)(f_1-cf_0)\,d\mu
                         +c(\alpha-E_0\psi)\geq0.
\]
Equality holds if and only if both of the following hold:
\[
\begin{gathered}
 \psi=1\quad\mu\hbox{-a.e. on }\{f_1>cf_0\},\qquad
 \psi=0\quad\mu\hbox{-a.e. on }\{f_1<cf_0\},\\
                              c(\alpha-E_0\psi)=0.
\end{gathered}
\]
Thus these conditions characterize every most-powerful test.
On the threshold equality set any measurable randomization is
permitted subject to the size constraint. If \(c>0\), every
most-powerful competitor has size \(\alpha\); if \(c=0\),
a smaller size can attain the same power.

At \(\alpha=0\), the maximum power is \(P_1(f_0=0)\);
a level-zero test attains it exactly when it is 1
\(P_1\)-a.e. on \(\{f_0=0\}\).
At \(\alpha=1\), the maximum power is 1 and is attained
exactly by tests equal to 1 \(P_1\)-a.e.
\end{hbtheorem}

For a nonrandomized continuous-density problem and \(k>0\),
the source's inequality \(p_0\leq kp_1\) is the same
likelihood ordering with \(c=1/k\). Its comparison of the
two regions outside their common part is the indicator
case of the integral proof above. Positive-measure
threshold sets permit the stated boundary choices
subject to size, so an optimal test need not be unique.

\begin{hblemma}{thm:p15-neyman-pearson-variation}{proved=annex:p15.neyman.pearson.variation}{The source's local boundary-variation argument}
Let \(d\geq2\) and let \(f_0,f_1\) be continuous densities on \(\mathbb R^d\),
\(k>0\), and \(h=f_0-kf_1\). On an open
\(B\subseteq\mathbb R^{d-1}\), let continuous functions
\(\ell<q\) describe a local region
\(\ell(z)<t<q(z)\).
For every continuous compactly supported \(\eta\) in \(B\),
vary its upper boundary to \(q(z)+s\eta(z)\).
The derivative at zero of the integral of \(h\) over
this varying region is
\[
                         \int_B\eta(z)h(z,q(z))\,dz.
\]
If every such variation is stationary, then
\(f_0(z,q(z))=k f_1(z,q(z))\) throughout the patch.
Reversing which side of the boundary is in the region
reverses the derivative's sign and gives the same condition.
The one-dimensional version is the ordinary variable-endpoint
integral formula. This is a local necessary condition under
the stated variation assumptions, not an assumption needed
by the general measurable test theorem.
\end{hblemma}

\section{Exact finite binomial tests}

\begin{hbdefinition}{def:p15-binomial-testing}{The binomial experiment}
Observe \(K\in\{0,\ldots,n\}\), \(n\geq1\), with
\[
                 b_{n,p}(k)=P_p(K=k)=\binom nkp^k(1-p)^{n-k},
                                  \qquad0\leq p\leq1.
\]
The endpoint laws put all mass at 0 for \(p=0\) and at
\(n\) for \(p=1\). Equivalently \(K\) counts successes
in \(n\) independent Bernoulli trials.
All formulas at endpoints use these polynomial limiting values.
For \(0<p_0<1\), choose \(j\in\{0,\ldots,n\}\) and
\(\gamma\in[0,1]\) with
\[
 P_{p_0}(K>j)\leq\alpha\leq P_{p_0}(K\geq j),\qquad
 \gamma=\frac{\alpha-P_{p_0}(K>j)}{P_{p_0}(K=j)},
\]
and put \(\varphi(k)=\mathbf1_{\{k>j\}}+\gamma\mathbf1_{\{k=j\}}\).
\end{hbdefinition}

\begin{hbtheorem}{thm:p15-binomial-most-powerful}{proved=annex:p15.binomial.most.powerful}{One-sided uniform optimality and an explicit two-sided failure}
For \(0<\alpha<1\), this test is most powerful for
\(p=p_0\) against every \(p_1>p_0\).
It is uniformly most powerful of level \(\alpha\)
for \(H_0:p\leq p_0\) against \(H_1:p>p_0\).
Its power is the exact finite sum
\[
                   P_p(K>j)+\gamma P_p(K=j).
\]
For \(n=4,p_0=1/2,\alpha=1/10\), the construction is
\[
 \varphi(4)=1,\quad\varphi(3)=3/20,\quad
 \varphi(0)=\varphi(1)=\varphi(2)=0.
\]
It has size \(1/10\), power \(243/640\) at \(p=3/4\),
and type-II error \(397/640\).
The nonrandomized rule rejecting only \(K=4\) has
size \(1/16\) and power \(81/256\) at \(3/4\).

Uniform optimality need not extend to a two-sided
alternative. For \(n=1,p_0=1/2,\alpha=1/2\), no test
is most powerful simultaneously against \(p=1/4\)
and \(p=3/4\).
\end{hbtheorem}

\begin{hbdefinition}{def:p15-valid-p-value}{A valid p-value}
A measurable \(\pi:\mathsf X\to[0,1]\) is valid for
\(\Theta_0\) if
\[
                  P_\theta(\pi(X)\leq t)\leq t
           \quad(\theta\in\Theta_0,\ 0\leq t\leq1).
\]
It yields a level-\(\alpha\) test by rejection at
\(\pi\leq\alpha\); strict rejection at \(\pi<\alpha\)
is also valid. The definition concerns sampling probabilities
under null laws, not a probability that the null is true.
\end{hbdefinition}

Nor is a small p-value a statement about the size of a departure from the null: how significance, confidence intervals and posterior probabilities relate to the practical magnitude of an effect is treated in \cref{sec:p15-w02b-p15-059-magnitude}.

\begin{hbtheorem}{thm:p15-discrete-p-values}{proved=annex:p15.discrete.p.values}{Discrete tails, exact randomization and conservatism}
For a finite ordered statistic \(T\) under a specified
null law, define for an observed value \(t\)
\[
 \pi(t)=P_0(T'\geq t),\qquad
 \pi_{\rm r}(t,u)=P_0(T'>t)+u P_0(T'=t),
\]
where \(T'\) has that null law and \(U\) is independent
uniform on \((0,1)\); evaluate these functions at \(T\) and \((T,U)\).
The inclusive tail is valid and can be strictly conservative.
The randomized tail is exactly uniform under the null.
The lower-tail versions have the same properties.

For binomial upper-tail testing at \(p_0\), the inclusive
tail remains valid under every \(p\leq p_0\).
In the preceding \(n=4\) example,
\(\pi(4)=1/16\), \(\pi(3)=5/16\), and
\(\pi_{\rm r}(3,U)=1/16+U/4\).
Rejecting the randomized p-value at \(1/10\) is precisely
the most-powerful randomized rule already given.
\end{hbtheorem}

\section{Acceptance regions and confidence sets}

\begin{hbdefinition}{def:p15-confidence-procedure}{Coverage, targets and measurable set membership}
Let the target parameter be a measurable map
\(\tau:\Theta\to\Lambda\), where these parameter spaces
are supplied with sigma-fields. A confidence procedure
\(C(x,u)\subseteq\Lambda\) has a measurable membership
graph in
\(\Lambda\times\mathsf X\times(0,1)\).
It has coverage at least \(1-\alpha\) if
\[
            (P_\theta\times\lambda_{01})\{\tau(\theta)\in C(X,U)\}
                                     \geq1-\alpha
                                  \quad(\theta\in\Theta).
\]
Thus any nuisance parameters in \(\theta\) are quantified
over. An unrandomized procedure omits \(u\).
Equality for all \(\theta\) is called exact coverage;
a finite-sample lower bound may instead be conservative.
Neyman's 1937 confidence coefficient is denoted here
by \(1-\alpha\); his use of \(\alpha\) for that coefficient
is distinct from this chapter's testing significance level.
\end{hbdefinition}

\begin{hbtheorem}{thm:p15-confidence-test-inversion}{proved=annex:p15.confidence.test.inversion}{Test inversion and its measurable converse}
Let \((t,x)\mapsto\varphi_t(x)\in[0,1]\), \(t\in\Lambda\),
be jointly measurable, and suppose
\[
                    E_\theta\varphi_{\tau(\theta)}(X)\leq\alpha
                                               \quad(\theta\in\Theta).
\]
For each attainable \(t\), this is a level-\(\alpha\) test
of \(H_t:\tau(\theta)=t\), uniformly over its nuisance parameters.
Then
\[
                  C(x,u)=\{t\in\Lambda:u>\varphi_t(x)\}
\]
has the stated coverage, with the exact identity
\[
 (P_\theta\times\lambda_{01})\{\tau(\theta)\in C(X,U)\}
                       =1-E_\theta\varphi_{\tau(\theta)}(X).
\]
Conversely any procedure with a measurable graph and
the stated coverage gives level-\(\alpha\) tests on
the augmented sample space by
\(\delta_t(x,u)=\mathbf1_{\{t\notin C(x,u)\}}\).
Integrating out \(u\) gives jointly measurable randomized
tests on the original sample space with the same rejection
probabilities. In the nonrandomized case
\[
                       x\in A_t
                            \quad\Longleftrightarrow\quad
                       t\in C(x)
\]
is the acceptance-region/coverage equivalence in both directions.
No connectedness or nonemptiness of the inverted sets
is needed for this general identity.
\end{hbtheorem}

\begin{hbtheorem}{thm:p15-neyman-interval-construction}{proved=annex:p15.neyman.interval.construction}{The full interval-section construction}
For a real target \(t\), suppose measurable acceptance
regions \(A_t\) satisfy \(P_\theta(A_{\tau(\theta)})=q\),
\(0<q<1\), for every full parameter \(\theta\).
Put \(S(x)=\{t:x\in A_t\}\).
Assume for every \(x\):
\begin{enumerate}
\item \(S(x)\) is nonempty.
\item If \(a,b\in S(x)\) and \(a<t<b\), then \(t\in S(x)\).
\item If all \(t\) with \(a<t<b\) belong to \(S(x)\), for
      finite \(a<b\), then \(a,b\in S(x)\).
\end{enumerate}
Then, with extended endpoints
\(\ell(x)=\inf S(x)\), \(u(x)=\sup S(x)\),
\[
                  S(x)=[\ell(x),u(x)]\cap\mathbb R.
\]
Assume also that for every rational \(r\) the two events
\[
 \{x:\exists t<r,\ x\in A_t\},\qquad
 \{x:\exists t>r,\ x\in A_t\}
\]
are measurable. Then \(\ell,u\) are measurable, and
\[
                      P_\theta\{\ell(X)\leq\tau(\theta)\leq u(X)\}=q.
\]
Conversely nonempty closed interval procedures with
measurable endpoints and coverage \(q\) have all the
displayed acceptance and section properties.
This supplies the membership equivalences and interval
reconstruction of Neyman's Propositions I--VII, including
singleton sections and infinite endpoints.
\end{hbtheorem}

\section{Exact binomial tail inversion}

\begin{hbdefinition}{def:p15-binomial-confidence-limits}{Closed equal-tail limits}
Let \(0<\alpha<1\), \(d=\alpha/2\), and \(n\geq1\).
Define
\[
 T_{n,p}(k)=\sum_{j=k}^n b_{n,p}(j),\qquad
 F_{n,p}(k)=\sum_{j=0}^k b_{n,p}(j).
\]
Set \(\ell_n(0)=0\), \(u_n(n)=1\). For the other
counts define the unique roots
\[
 T_{n,\ell_n(k)}(k)=d\quad(1\leq k\leq n),\qquad
 F_{n,u_n(k)}(k)=d\quad(0\leq k<n).
\]
The reported interval is
\(C_n(k)=[\ell_n(k),u_n(k)]\subseteq[0,1]\).
\end{hbdefinition}

\begin{hblemma}{thm:p15-binomial-tail-roots}{proved=annex:p15.binomial.tail.roots}{Existence, uniqueness and every boundary}
The roots above exist in \((0,1)\). For \(0<p<1\),
\[
\begin{aligned}
 \frac{d}{dp}T_{n,p}(k)
     &=n\binom{n-1}{k-1}p^{k-1}(1-p)^{n-k}>0
                                      &&(1\leq k\leq n),\\
 \frac{d}{dp}F_{n,p}(k)
     &=-n\binom{n-1}{k}p^k(1-p)^{n-k-1}<0
                                      &&(0\leq k<n).
\end{aligned}
\]
Each \(C_n(k)\) is a nonempty closed interval.
In particular
\[
                u_n(0)=1-d^{1/n},\qquad \ell_n(n)=d^{1/n}.
\]
\end{hblemma}

\begin{hbtheorem}{thm:p15-binomial-confidence-coverage}{proved=annex:p15.binomial.confidence.coverage}{The exact inversion and finite-sample guarantee}
For every \(p\in[0,1]\),
\[
                            P_p\{p\in C_n(K)\}\geq1-\alpha.
\]
The interval is exactly the inversion of the nonrandomized
test rejecting \(p\) when
\[
                       T_{n,p}(K)<d\quad\hbox{or}\quad F_{n,p}(K)<d.
\]
The strict inequalities make the root endpoints accepted.
At \(p=0\) and \(p=1\), coverage is 1.
The guarantee uses the exact binomial law; it does not
assert equality to \(1-\alpha\) at every parameter.

At \(n=4,\alpha=1/10\),
\[
 C_4(0)=[0,1-20^{-1/4}],\qquad C_4(4)=[20^{-1/4},1].
\]
At \(n=2,\alpha=1/10\),
\[
                       C_2(1)=[1-\sqrt{19/20},\sqrt{19/20}].
\]
All three possible intervals for \(n=2\) contain \(p=1/2\),
so its coverage there is 1, strictly above \(9/10\).
\end{hbtheorem}

\begin{hbtheorem}{thm:p15-binomial-beta-quantiles}{proved=annex:p15.binomial.beta.quantiles}{The beta-quantile formulas without zero shape parameters}
Let \(Q_{a,b}(t)\) be the unique \(t\)-quantile of the
\(\operatorname{beta}(a,b)\) distribution from
\cref{def:p15-beta-integral}, for \(a,b>0\) and \(0<t<1\).
For integer binomial counts,
\[
 \ell_n(k)=Q_{k,n-k+1}(d)\quad(k\geq1),\qquad
 u_n(k)=Q_{k+1,n-k}(1-d)\quad(k<n).
\]
The separately defined endpoints \(\ell_n(0)=0\) and
\(u_n(n)=1\) remain in force; a beta distribution with
a zero shape is not used.
\end{hbtheorem}

\begin{hbtheorem}{thm:p15-confidence-credibility-boundary}{proved=annex:p15.confidence.credibility.boundary}{Sampling coverage and posterior content are different assertions}
Under independent repetitions at a fixed parameter,
the frequency of coverage has the coverage probability
as its limit in probability. After fixing the observed
data, the parameter in the sampling experiment remains
fixed; the coverage guarantee is not a posterior probability.

For a separate joint model, give the binomial parameter
a uniform prior on \((0,1)\). Given \(K=k\), its conditional
law is \(\operatorname{beta}(k+1,n-k+1)\).
For \(n=4,k=0,\alpha=1/10\), the confidence interval
\(C_4(0)\) has posterior content
\[
                              1-20^{-5/4},
\]
which is not \(9/10\).
Conversely the central posterior \(9/10\) beta interval
at each count has posterior content \(9/10\), but its
sampling coverage at \(p=0\) is zero, because its lower
endpoint after \(K=0\) is strictly positive.
Indeed the positive minimum of the finitely many lower
posterior endpoints also makes coverage zero at some
strictly positive parameter values.
\end{hbtheorem}
